\documentclass{article}
\usepackage{../math_template}

\author{Jonathan Kasongo}
\title{British Math Olympiad Round 1 2002 --- P4/5}

\begin{document}
\maketitle

\begin{problem}{British Math Olympiad Round 1 2002 --- P4/5}
Twelve people are seated around a circular table. In how many ways
can six pairs of people engage in handshakes so that no arms cross?
(Nobody is allowed to shake hands with more than one person at once.)
\end{problem}

\begin{solution}{Solution}
We will solve the following more general problem: Let $n$ be a positive
integer. Given $2n$ people are sat around a circular table, in how many
ways can $n$ pairs of people engage in handshakes so that no arms cross? \\

Let's denote the $2n$ people sat at the table by $a_1, a_2, ..., a_{2n}$,
where each of the people are sat in that order. \\

\textit{Lemma: If person $a_k$ handshakes with person $a_{k + d}$ then
each of the people $a_{k+1}, a_{k+2}, ..., a_{k + d - 1}$ must handshake
another person amongst themselves, in order to have no arms crossing.} \\

\textit{Proof:} When person $a_k$ handshakes with person $a_{k + d}$ their
linked arms make a chord across the circular table, splitting the table
into 2 halves. One halve contains the people $a_{k+1}, a_{k+2}, ...,
a_{k + d - 1}$, and the other half contains the remainding people. If
2 people from different halves were to handshake then their arms would
cross that chord, which is a contradiction. $\square$ \\

Due to our lemma if person $a_k$ handshakes with person $a_{k+d}$, it
creates a group of $d-1$ people who have to handshake amongst themselves.
For this to be possible $d-1$ must be even and $d$ must be odd. Thus
if person $a_p$ handshakes with person $a_q$, the integers
$p$ and $q$ must have opposite parity. \\

Let $H_{2k}$ be the number of ways $k>2$ pairs of people can handshake such
that no arms cross, given that the $2k$ people are sat around a circular
table. Consider person $a_1$, they must handshake with person $a_{2m}$,
for integer $1 \leq m \leq k$. Then there are 2 groups of people that
should handshake among themselves under the same constraints as before.
One group is of size $2m - 2$ and the other is of size $2k - 2m$.
Thus

\[
\begin{aligned}
H_{2k} &= \sum_{m=1}^k H_{2m - 2} \cdot H_{2k - 2m} \\
\end{aligned}
\]

It is easily verifiable that $H_0 = 1, H_2 = 1$ and from there one can
show that $H_12 = 132$ by brute-force computation (or by using the fact
that $H_{2n}$ is the $n^\text{th}$ catalan number $
\frac{1}{n+1} \binom{2n}{n}$). $\blacksquare$

\end{solution}

\end{document}
